% from ega4-15-fr.tex
\oldpage[IV-3]{225}
Posons $S=\operatorname{Spec}(A)$, et considérons le $S$-schéma $P$ obtenu en
faisant éclater le point fermé $s$ de $S$ (\textbf{II}, 8.1.3)~; si $\mathfrak m$
est l'idéal maximal de $A$, on a donc par définition
$P=\operatorname{Proj}(B)$, où $B$ est la $A$-algèbre graduée
$\bigoplus_{n\geqslant0}\mathfrak m^n$. Si $h:P\to S$ est le morphisme
structural, la fibre $h^{-1}(s)$ est isomorphe à
$\operatorname{Proj}(B\otimes_A k)$ (\textbf{II}, 2.8.10), et par définition
$B\otimes_A k$ est la $k$-algèbre graduée
$\bigoplus_{n\geqslant0}\mathfrak m^n/\mathfrak m^{n+1}$~; mais comme $A$ est
régulier, cette algèbre est isomorphe à une algèbre de polynômes
$k[t_1,\ldots,t_e]$ ($t_i$ indéterminées) ($\mathbf{0}_{\mathrm{IV}}$, 17.1.1)
et par suite $h^{-1}(s)$ est isomorphe à $\mathbf P_k^{e-1}$. Soit $\zeta$ le
point générique de $h^{-1}(s)$~; si $t_i'$ ($1\leqslant i\leqslant e$) sont les
éléments de $\mathfrak m$ dont les classes mod. $\mathfrak m^2$ sont les $t_i$,
$B_{(t_e')}$ est l'anneau d'un voisinage ouvert affine de $\zeta$ dans $P$, et
l'on a $\mathbf{k}(\zeta)=k(t_1,\ldots,t_{e-1})$~; d'autre part, comme $A$ est
intégralement clos, on vérifie aisément qu'il en est de même de $B$ (Bourbaki,
\emph{Alg. comm.}, chap.~V, \S~1, n\textsuperscript{o}~8, cor.~1 de la prop.~21),
donc $B_{(t_e')}$ est aussi intégralement clos, et il en est par suite de même de
l'anneau local $\mathcal O_{P,\zeta}$. Enfin, comme $A=\mathcal O_s$ est
régulier, donc un anneau universellement caténaire (5.6.4), on a, en vertu de
(5.6.5),
$\dim(\mathcal O_{P,\zeta})=\dim(\mathcal O_s)-(e-1)$, puisque $h$ est un
morphisme birationnel, que $\dim(h^{-1}(s))=e-1$ et que l'anneau local de la
fibre $h^{-1}(s)$ en son point générique $\zeta$ est de dimension $0$. Mais
$\dim(\mathcal O_s)=e$, donc $\dim(\mathcal O_{P,\zeta})=1$, et
$\mathcal O_{P,\zeta}=V$ est un anneau de valuation discrète, étant noethérien,
de dimension $1$ et intégralement clos (\textbf{II}, 7.1.6)~; il répond
évidemment à la question.

Ce lemme étant établi, on peut reprendre le raisonnement fait pour l'inégalité
(15.1.1.1) avec $Y'=\operatorname{Spec}(V)$~; cette fois $\mathbf{k}(y')$ est
une extension \emph{séparable} de $\mathbf{k}(y)$, et par suite (4.7.9), on a
$\operatorname{long}((\mathcal F_y)_z)=
\operatorname{long}((\mathcal F'_{y'})_{z'})$~; comme par ailleurs
$f'^{-1}(\eta')=f^{-1}(\eta)$ et $\mathcal F'_{\eta'}=\mathcal F_\eta$, on est
encore ramené à prouver l'inégalité (15.1.1.5). Or, si $t$ est une uniformisante
de $\mathcal O_{y'}$, $f'^{-1}(y')$ est le sous-préschéma fermé de $X'$ défini
par l'idéal $t\mathcal O_{X'}$, et on a
$\mathcal F'_{y'}=\mathcal F'/t\mathcal F'$~; d'autre part, on vérifie aussitôt
(\textbf{I}, 9.1.12) que
$(\mathcal F'_{\eta'})_{z_i'}=\mathcal F'_{z_i'}$~; l'inégalité à démontrer
(15.1.1.5) n'est autre alors qu'un cas particulier de (3.4.1.1).
\end{proof}

\begin{corollary}[15.1.2]
\label{IV.15.1.2-fr}
Les hypothèses étant celles de (15.1.1), supposons de plus que
$\lambda_z(\mathcal F_y)=1$ (resp. que $\mathcal O_y$ soit régulier et
$\operatorname{long}((\mathcal F_y)_z)=1$). Alors il existe au plus une composante
irréductible $X_i$ de $\operatorname{Supp}(\mathcal F)$ contenant $z$ et telle que
la restriction $X_i\to Y$ de $f$ soit universellement ouverte au point $z$. De
plus, si $z_i$ est le point générique de cette composante, on a
$\lambda_{z_i}(\mathcal F_\eta)=1$ (resp.
$\operatorname{long}((\mathcal F_\eta)_{z_i})=1$).
\end{corollary}

Cela résulte aussitôt de (15.1.1) en notant que l'on a nécessairement, par
définition de $X_i$, $(\mathcal F_\eta)_{z_i}\ne0$ (\textbf{I}, 9.1.13).

\begin{corollary}[15.1.3]
\label{IV.15.1.3-fr}
Soient $Y$ un préschéma irréductible localement noethérien de point générique
$\eta$, $f:X\to Y$ un morphisme localement de type fini, $\mathcal F$ un
$\mathcal O_X$-Module cohérent de support $X$, $x$ un point de $X$, $y=f(x)$. On
suppose que~: $1^\circ$ $x$ n'appartient qu'à une seule composante irréductible
$Z$ de $f^{-1}(y)$~; $2^\circ$
$\mathcal F_y=\mathcal F\otimes_{\mathcal O_Y}\mathbf{k}(y)$ est géométriquement
réduit sur $\mathbf{k}(y)$, autrement dit, si $z$ est le point générique de $Z$,
$\lambda_z(\mathcal F_y)=1$ (resp. $\mathcal O_y$ est régulier et
$\operatorname{long}((\mathcal F_y)_z)=1$). Alors, il existe au plus une
composante irréductible $X'$ de $X$ contenant $x$, telle que
$\dim_x(X'\cap f^{-1}(y))=\dim_x(f^{-1}(y))$ et que la restriction $X'\to Y$ de
$f$ soit universellement ouverte aux points génériques des composantes
irréductibles de $X'\cap f^{-1}(y)$ contenant $x$. De plus, s'il existe une telle
composante $X'$ et si $z'$ est son point générique, on a
$\lambda_{z'}(\mathcal F_\eta)=1$ (resp.
$\operatorname{long}((\mathcal F_\eta)_{z'})=1$).
\end{corollary}

